% ===== SLIDE 7 — MÉTRICAS =====
\begin{frame}{Métricas de Evaluación}
\framesubtitle{¿Cómo sabemos si un modelo es confiable?}

La mala jugada de la \alert{precisión}. Estas son las metricas:

\vspace{4pt}

\begin{columns}[T]
\begin{column}{0.45\textwidth}
\[
\text{Sensibilidad} = \frac{\text{VP}}{\text{VP} + \text{FN}}
\]
\begin{center}
{\small\color{muted} Capacidad de detectar enfermos}
\end{center}
\end{column}
\begin{column}{0.45\textwidth}
\[
\text{Especificidad} = \frac{\text{VN}}{\text{VN} + \text{FP}}
\]
\begin{center}
{\small\color{muted} Capacidad de identificar sanos}
\end{center}
\end{column}
\end{columns}

\vfill

\begin{alertblock}{\faChartArea\; Curva ROC y AUC}
El AUC resume el rendimiento en un valor de 0 a 1. Un AUC de \textbf{1.0} = perfección.
\end{alertblock}

\end{frame}
